\begin{abstract}
A satellite or a drone can carry only a small network, and the land-cover classes it meets next come with a handful of labelled examples. We ask whether the choice of channels matters when a pretrained ResNet-18 is pruned for that job. Eight channel scores are compared at $50\%$ to $10\%$ of the parameters: random selection, filter $\ell_1$, BatchNorm scale, a Taylor term, three single statistics, and DACIS, a rank sum of those statistics from an earlier preprint. They are trained under seven schedules, among them prune, meta-train, prune again (PMP). Test classes are held out: $15$ aerial-scene classes of NWPU-RESISC45, five land-cover classes of 13-band Sentinel-2 EuroSAT-MS and eight PlantVillage crop-disease classes, with PlantDoc field photographs as a domain-shift check; five class partitions, one set of test episodes for all methods. At $30\%$ of the parameters, \num{cmp:above:PMP} of \num{cmp:total:PMP} paired comparisons of a computed score with random selection favour the computed score and \num{cmp:below:PMP} favour random selection; the best computed score is ahead by at most \num{cmp:maxgain:PMP} points, and the ranking of the scores agrees across partitions little more than unrelated rankings would (Kendall's $W$ \num{kw:min} to \num{kw:max}). Episodic training did not raise accuracy. Pruning to $30\%$ removed \num{cost:macs-red:30}\% of the multiply-accumulates but only \num{cost:cpu-red:30}\% of single-thread CPU latency. A notebook reruns every number.
\end{abstract}
